% ===================================================================
%  TABLEAU N° 4 — Age mûr et Vieillesse.
%  Feuillets 25 a 30 du fac-simile = folios 23 a 28.
%  Correspondance : folio imprime = numero de feuillet - 2.
%
%  Ce tableau est, comme son homologue ido (texto/io/13-tabelo-04.tex),
%  partage en DEUX SCENES (Première scène, Deuxième scène) dont
%  l'auteur remet la numerotation des alineas a 1. Les cles %%K
%  suivent la consigne : t04-01-1 apparait donc DEUX FOIS dans ce
%  fichier, une fois par scene ; outils/ceni.py inserera l'indice de
%  scene apres coup.
%
%  ECART EDITORIAL REEL : l'ido intercale quatre sous-titres internes
%  a l'interieur des deux scenes (« L'Autuno. », « La Parentaro. »,
%  « La Amiki. », « La Desero. » pour la premiere scene ; « La
%  Vizito. », « La Maladeso. », « La Avulo. » pour la seconde). Le
%  francais ne connait qu'un seul intitule par scene (« Le repas de
%  noce. » et « L'anniversaire du grand-père. ») : les alineas
%  s'enchainent ensuite sans rubrique. Ce n'est pas une lacune du
%  releve : Rochelle ne decoupe pas la scene ou Guignon la decoupe.
%
%  ATTENTION — LE FAC-SIMILE FRANCAIS EST UNE PRISE DE VUE, non un
%  passage au scanner : le contraste est inegal et le feuillet 30 est
%  particulierement bruite. Les mots ont ete verifies en agrandissant
%  l'image plutot que devines.
% ===================================================================

% --- Feuillet 25 (folio 23) -----------------------------------------
\begin{VUpage}[25]{23}
%%K t04-apar-1 apar
\VUsaut{4.0mm}
\VUtitre{15.0pt}{60}{TABLEAU N\textsuperscript{o} 4}
\VUsaut{1.6mm}
\VUtitre{11.4pt}{0}{Age mûr et Vieillesse.}
\VUsaut{3.2mm}

%%K t04-tit-1 sub
\VUcentre{11.4pt}{0}{\textit{Première scène.}}
\VUsaut{1.6mm}
\VUpk{10.2pt}{40}{Le repas de noce.}
\VUsaut{3.0mm}

%%K t04-c1-01-1 p
\VUblancAlinea
1. --- Les jeunes gens ont grandi. L'un d'eux, Jean,\nl
se marie. Nous assistons au repas de \VUgras{noce}, qui ras\cc
semble autour de la même table les familles des deux\nl
époux.

%%K t04-c1-02-1 p
\VUblancAlinea
2. --- Nous sommes en \VUgras{automne}, au mois de \VUgras{sep}\cc
\VUgras{tembre}. Le vent froid d'\VUgras{octobre} et de \VUgras{novembre}\nl
n'a pas encore dépouillé la campagne de sa verte\nl
parure. L'air du soir est tiède et parfumé; aussi le\nl
dîner a-t-il lieu sous une \VUgras{véranda} \textsuperscript{(8)}, donnant sur le\nl
jardin.

%%K t04-c1-03-1 p
\VUblancAlinea
3. --- Au loin, sur la \VUgras{colline} \textsuperscript{(3)}, est située la\nl
maison d'habitation des parents de Jean, un élégant\nl
\VUgras{château} \textsuperscript{(1)} caché dans la verdure; de beaux vigno\cc
bles, qui produisent d'excellent vin, s'étendent sur le\nl
versant du coteau. Le vigneron les cultive et les\nl
soigne.

%%K t04-c1-04-1 p
\VUblancAlinea
4. --- Au-dessus des vignes, passe un vol de \VUgras{per}\cc
\VUgras{drix} \textsuperscript{(2)}, effarouchées par l'approche du \VUgras{garde}\cc
\VUgras{chasse} \textsuperscript{(5)}. Celui-ci, le \VUgras{fusil} sous le bras, la \VUgras{carnas}\cc
\VUgras{sière} en bandoulière, bat les \VUgras{buissons} \textsuperscript{(4)} avec son\nl
\VUgras{chien} \textsuperscript{(6)}. Il surveille la propriété et empêche les\nl
braconniers de détruire le gibier.

%%K t04-c1-05-1 p
\VUblancAlinea
5. --- A l'extérieur de la \VUgras{véranda}, grimpe une\nl
\VUgras{treille}, à laquelle sont suspendues de lourdes\nl
\VUgras{grappes} \textsuperscript{(7)} de raisin.

%%K t04-c1-06-1 p
\VUblancAlinea
6. --- Les jolies fleurs d'automne qui garnissent la\nl
\VUgras{table} \textsuperscript{(16)} sont des \VUgras{chrysanthèmes} \textsuperscript{(9)}; le \VUgras{père} \textsuperscript{(10)}
\parplein
\end{VUpage}

% --- Feuillet 26 (folio 24) -----------------------------------------
\begin{VUpage}[26]{24}
%%K t04-c1-06-1 p suite
\VUcontinue
de la mariée en a mis un à la boutonnière de son\nl
habit.

%%K t04-c1-07-1 p
\VUblancAlinea
7. --- Le repas touche à sa fin. Le \VUgras{domestique} \textsuperscript{(11)}\nl
commence à apporter les desserts et va enlever les\nl
différentes parties du couvert, \VUgras{cuillères} \textsuperscript{(14)}, \VUgras{four}\cc
\VUgras{chettes} \textsuperscript{(12)}, \VUgras{couteaux} \textsuperscript{(13)}, \VUgras{plats} \textsuperscript{(15)}, dont on n'a\nl
plus besoin.

%%K t04-c1-08-1 p
\VUblancAlinea
8. --- Tout le monde semble joyeux, et le sourire\nl
avec lequel la \VUgras{mère} \textsuperscript{(17)} du marié regarde ses enfants\nl
prouve son bonheur.

%%K t04-c1-09-1 p
\VUblancAlinea
9. --- Ce mariage unit deux riches familles du\nl
pays. Jean, le \VUgras{mari} \textsuperscript{(18)}, est le \VUgras{fils} d'un gros pro\cc
priétaire et il devient le \VUgras{gendre} d'un grand indus\cc
triel.

%%K t04-c1-10-1 p
\VUblancAlinea
10. --- Les \VUgras{époux} sont jeunes et pourtant ils étaient\nl
\VUgras{fiancés} depuis longtemps. La \VUgras{jeune femme} \textsuperscript{(19)},\nl
Marthe, est ravissante dans sa \VUgras{robe blanche} garnie\nl
de fleurs d'oranger.

%%K t04-c1-11-1 p
\VUblancAlinea
11. --- Son \VUgras{beau-père} \textsuperscript{(20)} est très heureux d'avoir\nl
une \VUgras{bru} aussi gentille, et la \VUgras{belle-mère} \textsuperscript{(21)} de Jean\nl
sourit en voyant sa \VUgras{fille} si jolie.

%%K t04-c1-12-1 p
\VUblancAlinea
12. --- Au bout de la table sont assis les autres\nl
\VUgras{invités} \textsuperscript{(22)}, des \VUgras{parents}, des \VUgras{amis}, des \VUgras{cousins},\nl
des \VUgras{cousines}. Un \VUgras{maître d'hôtel} \textsuperscript{(23)}, sa serviette\nl
sous le bras, les sert avec empressement.

%%K t04-c1-13-1 p
\VUblancAlinea
13. --- Du côté droit de la table, voici une \VUgras{demoi}\cc
\VUgras{selle} \textsuperscript{(24)}, \VUgras{amie} de la jeune femme, puis le \VUgras{frère}\nl
de celle-ci qui est son \VUgras{garçon d'honneur} \textsuperscript{(25)}. Il\nl
bavarde avec sa \VUgras{demoiselle d'honneur} \textsuperscript{(26)} qui, par\nl
ce mariage, devient la \VUgras{belle-sœur} de Marthe. C'est\nl
une jeune fille charmante. Elle a de jolis \VUgras{bijoux},\nl
un \VUgras{collier} de perles et un \VUgras{bracelet} en or.

%%K t04-c1-14-1 p
\VUblancAlinea
14. --- Près d'eux, le monsieur qui se tient debout\nl
et lève son verre est un \VUgras{ami} \textsuperscript{(27)} de la famille. Il est
\parplein
\end{VUpage}

% --- Feuillet 27 (folio 25) -----------------------------------------
\begin{VUpage}[27]{25}
%%K t04-c1-14-1 p suite
\VUcontinue
l'un des \VUgras{témoins} du marié. Il porte un \VUgras{toast}, boit\nl
à la \VUgras{santé} des jeunes époux et leur souhaite beau\cc
coup de bonheur. Il est en costume de cérémonie, en\nl
habit avec des \VUgras{souliers vernis} \textsuperscript{(28)}. Il est le cavalier\nl
de la \VUgras{grand'mère} \textsuperscript{(29)}, qui est \VUgras{veuve}.

%%K t04-c1-15-1 p
\VUblancAlinea
15. --- Le jeune homme en \VUgras{habit} \textsuperscript{(31)} à la fine\nl
moustache noire est le \VUgras{beau-frère} \textsuperscript{(30)} de Jean.\nl
C'est un ingénieur distingué. Il est le voisin d'une\nl
\VUgras{jeune femme} \textsuperscript{(33)} très élégante, la \VUgras{sœur aînée} de\nl
Marthe et la mère de la mignonne petite fille qui\nl
vient, en donnant le bras à son cousin, lui apporter\nl
une fleur et lui \VUgras{souhaiter le bonsoir}. Cette dame\nl
a de fort belles \VUgras{boucles d'oreilles} et une jolie\nl
\VUgras{coiffure}.

%%K t04-c1-16-1 p
\VUblancAlinea
16. --- Le petit cousin, si drôle avec sa \VUgras{blouse} \textsuperscript{(36)}\nl
et sa \VUgras{culotte} \textsuperscript{(37)} bouffante, est bien à plaindre. Il est\nl
\VUgras{orphelin} \textsuperscript{(35)}. Il a perdu tout jeune son père et sa\nl
mère. Son oncle lui sert de \VUgras{tuteur} \textsuperscript{(38)}. Il est resté\nl
célibataire pour élever le pauvre petit. C'est un\nl
excellent homme. Il est un peu chauve et très myope;\nl
il porte un \VUgras{lorgnon}.

%%K t04-c1-17-1 p
\VUblancAlinea
17. --- Derrière les convives, sur une table recou\cc
verte d'une \VUgras{nappe}, le \VUgras{dessert} \textsuperscript{(39)} est préparé. Dans\nl
une coupe, voici des \VUgras{fruits}, des \VUgras{pêches} \textsuperscript{(40)}, des\nl
\VUgras{poires} \textsuperscript{(41)}, des \VUgras{pommes} \textsuperscript{(42)}, des \VUgras{abricots} \textsuperscript{(43)} et\nl
des \VUgras{raisins} \textsuperscript{(44)}. A côté, des \VUgras{ananas} \textsuperscript{(45)}, un beau\nl
\VUgras{gâteau} \textsuperscript{(46)}, des bouteilles de \VUgras{liqueur} \textsuperscript{(48)} et du\nl
\VUgras{champagne} \textsuperscript{(49)} ainsi que des piles d'\VUgras{assiettes} \textsuperscript{(47)}\nl
propres.

%%K t04-c1-18-1 p
\VUblancAlinea
18. --- Le \VUgras{sommelier} \textsuperscript{(50)} débouche une \VUgras{bou}\cc
\VUgras{teille} \textsuperscript{(52)} de vin vieux avec un \VUgras{tire-bouchon} \textsuperscript{(51)}.\nl
Le \VUgras{bouchon} \textsuperscript{(53)} semble très difficile à tirer. Ce \VUgras{som}\cc
\VUgras{melier} a une \VUgras{serviette} \textsuperscript{(54)} sous le bras.

%%K t04-c1-19-1 p
\VUblancAlinea
19. --- Après le dîner, les messieurs passeront au\nl
fumoir et les dames iront s'apprêter pour le bal.
\end{VUpage}

% --- Feuillet 28 (folio 26) -----------------------------------------
\begin{VUpage}[28]{26}
%%K t04-tit-2 sub
\VUsaut{4.0mm}
\VUcentre{11.4pt}{0}{\textit{Deuxième scène.}}
\VUsaut{1.6mm}
\VUpk{10.2pt}{40}{L'anniversaire du grand-père.}
\VUsaut{3.0mm}

%%K t04-c2-01-1 p
\VUblancAlinea
1. --- Dix années se sont écoulées. Le père et la\nl
mère de Jean sont devenus des vieillards qui aiment\nl
beaucoup leurs trois petits enfants, deux \VUgras{petits}\cc
\VUgras{fils} \textsuperscript{(2)} et une \VUgras{petite-fille} \textsuperscript{(3)}. Comme c'est aujourd'hui\nl
l'\VUgras{anniversaire} du \VUgras{grand-père}, les enfants vien\cc
nent lui souhaiter une bonne fête.

%%K t04-c2-02-1 p
\VUblancAlinea
2. --- Le petit Pierre est encore tout jeune, il n'a que\nl
trois ans. Il aperçoit le \VUgras{chat} \textsuperscript{(1)} qui s'avance. Il\nl
veut caresser son beau \VUgras{poil} soyeux et sa longue\nl
\VUgras{queue}, il ne pense pas que, sous les jolies \VUgras{pattes},\nl
se cachent de vilaines griffes pointues qui pourront\nl
bien l'égratigner.

%%K t04-c2-03-1 p
\VUblancAlinea
3. --- Sa sœur Gabrielle, qui lui donne la main, est\nl
beaucoup plus sérieuse, et Marcel, avec son grand\nl
\VUgras{manteau} \textsuperscript{(4)} et sa \VUgras{ceinture} \textsuperscript{(5)} de cuir, a l'air d'un\nl
petit homme, malgré ses culottes courtes qui laissent\nl
apercevoir ses \VUgras{bas} \textsuperscript{(6)}.

%%K t04-c2-04-1 p
\VUblancAlinea
4. --- Leur père a mis son gros \VUgras{pardessus} \textsuperscript{(7)}\nl
d'hiver à col de \VUgras{fourrure} \textsuperscript{(8)} et, dans sa \VUgras{poche} \textsuperscript{(9)}, il\nl
a un foulard. Sa femme a son chapeau de feutre\nl
garni de \VUgras{plumes} \textsuperscript{(10)}, sa \VUgras{jaquette} \textsuperscript{(11)}, son \VUgras{boa} \textsuperscript{(12)} et\nl
son \VUgras{manchon} \textsuperscript{(13)}.

%%K t04-c2-05-1 p
\VUblancAlinea
5. --- Il fait très froid, car nous sommes en hiver.\nl
La \VUgras{neige} \textsuperscript{(19)} est tombée toute la journée et les toits\nl
des maisons sont tout blancs. Avec la \VUgras{nuit}, le\nl
froid se fait encore plus vif. Au ciel, la \VUgras{lune} \textsuperscript{(17)} et\nl
les \VUgras{étoiles} \textsuperscript{(18)} brillent et percent l'\VUgras{obscurité}.

%%K t04-c2-06-1 p
\VUblancAlinea
6. --- Le pauvre \VUgras{aveugle} \textsuperscript{(20)} qui traverse la place\nl
devant la \VUgras{pharmacie} \textsuperscript{(16)}, conduit par son \VUgras{caniche} \textsuperscript{(21)},\nl
est bien malheureux. C'est du moins ce que pense\nl
Justine, la \VUgras{bonne} \textsuperscript{(14)}, qui apporte de la cuisine un\nl
\VUgras{bol} \textsuperscript{(15)} de \VUgras{tisane} bien chaude et bien sucrée.
\end{VUpage}

% --- Feuillet 29 (folio 27) -----------------------------------------
\begin{VUpage}[29]{27}
%%K t04-c2-07-1 p
\VUblancAlinea
7. --- \VUgras{Grand'maman} \textsuperscript{(23)}, qui veut aller chercher\nl
ses \VUgras{lunettes} \textsuperscript{(26)} sur la table pour lire l'\VUgras{ordon}\cc
\VUgras{nance} \textsuperscript{(27)} que le \VUgras{docteur} \textsuperscript{(25)} vient d'écrire, se lève\nl
de son fauteuil en s'appuyant sur sa \VUgras{canne} \textsuperscript{(24)}.

%%K t04-c2-08-1 p
\VUblancAlinea
8. --- Elle a souvent de légères \VUgras{indispositions},\nl
\VUgras{vertiges, rhumes, maux de dents}; ses jambes\nl
sont faibles et ses yeux ne sont plus très bons. Malgré\nl
cela, elle se moque volontiers de son vieux \VUgras{médecin},\nl
qui veut toujours lui prescrire une \VUgras{potion} \textsuperscript{(28)}.\nl
Aujourd'hui elle est toute joyeuse d'avoir sa petite\nl
famille autour d'elle.

%%K t04-c2-09-1 p
\VUblancAlinea
9. --- Il fait bon dans la chambre bien close. Dans la\nl
\VUgras{cheminée} \textsuperscript{(29)} de marbre flambe un \VUgras{feu} \textsuperscript{(30)} superbe,\nl
et l'on se sent heureux sous la douce \VUgras{lumière} de la\nl
\VUgras{lampe} \textsuperscript{(31)} que l'on vient d'allumer.

%%K t04-c2-10-1 p
\VUblancAlinea
10. --- Le \VUgras{grand-père} \textsuperscript{(33)} lui-même oublie qu'il a\nl
été malade, que son \VUgras{rhumatisme} le tient cloué dans\nl
son \VUgras{fauteuil} \textsuperscript{(34)} et que, hier encore, il avait la\nl
\VUgras{fièvre} et des \VUgras{syncopes}. Il ne pense plus que\nl
l'hiver dernier il a eu une \VUgras{fluxion de poitrine} et\nl
qu'une \VUgras{toux} rauque et pénible le faisait beaucoup\nl
souffrir.

%%K t04-c2-10-2 p
Et pourtant il ne s'est rétabli que difficilement.\nl
Maintenant il se porte un peu mieux. Mais sa jambe,\nl
qu'il tient appuyée sur un \VUgras{coussin} \textsuperscript{(38)} posé sur un\nl
petit \VUgras{tabouret} \textsuperscript{(39)}, lui fait encore bien mal.

%%K t04-c2-11-1 p
\VUblancAlinea
11. --- Lorsqu'il est bien enveloppé dans sa chaude\nl
\VUgras{robe de chambre} \textsuperscript{(35)} à gros \VUgras{boutons} \textsuperscript{(36)} de nacre\nl
et qu'il a auprès de lui, pour lui lire son journal,\nl
la petite \VUgras{orpheline} \textsuperscript{(40)}, au bonnet blanc, à la\nl
\VUgras{robe} \textsuperscript{(42)} et à la \VUgras{pèlerine} \textsuperscript{(41)} bleues, il ne se\nl
plaint pas.

%%K t04-c2-12-1 p
\VUblancAlinea
12. --- Chaque année, lorsque le \VUgras{calendrier} \textsuperscript{(43)}\nl
ramène la \VUgras{date} du 1\textsuperscript{er} décembre, il attend ses\nl
\VUgras{petits-enfants} avec impatience, car il sait qu'ils\nl
n'oublieront pas cette \VUgras{date} importante.
\end{VUpage}

% --- Feuillet 30 (folio 28) -----------------------------------------
\begin{VUpage}[30]{28}
%%K t04-c2-13-1 p
\VUblancAlinea
13. --- Il se rappelle avec émotion le temps bien\nl
lointain où il venait lui-même souhaiter une bonne\nl
fête au brillant \VUgras{maréchal} de l'Empire, dont le portrait\nl
est suspendu au mur, et qui est le \VUgras{bisaïeul} des\nl
enfants.

\VUsaut{6.0mm}
\VUfilet{22mm}
\end{VUpage}
